% !TeX root = ../深度学习与大语言模型技术原理.tex
% 由主文档通过 \input 载入；请编译主文档。
\setcounter{chapter}{9} % 保留原书章号。
\chapter{监督微调与参数高效适配}
\label{ch:sft}

\section{监督目标与回答区域掩码}

\emph{监督微调（supervised fine-tuning, SFT）}用指令--回答或多轮对话示例训练条件生成。设一条序列为 $z_{1:n}$，$m_t\in\{0,1\}$ 指示该位置是否属于需要学习的目标，则批目标可写为
\begin{equation}
\mathcal L_{\mathrm{SFT}}=-\frac{\sum_{i,t}m_{it}\log p_\theta(z_{it}\mid z_{i,<t})}{\sum_{i,t}m_{it}}.
\label{eq:sft}
\end{equation}
只训练助手回答是常见选择；用户消息虽然不计损失，仍进入模型作为条件，因此不计损失与不可注意这两种掩码含义不同。

多轮对话可以训练每轮助手响应，也可以只训练末轮，目标随之改变。应根据序列化后的角色边界构造掩码，避免用脆弱的普通字符串搜索误把用户引用的角色标记当成真实边界。工具调用内容若属于模型要生成的动作，也可进入监督目标；工具执行结果通常作为后续条件。结束词元是否计入损失会影响停止行为，应与服务端终止规则一致。

\section{拼接、截断与灾难性遗忘}

样本拼接（packing）把多条短样本放入同一计算序列，减少填充浪费。如果它们应相互独立，应使用分块因果掩码，使每条样本只访问自身前缀；位置编号是否重置也需与架构约定一致。如果允许跨文档注意力，则训练的是一个文档串接目标，应明确承认这种变化。仅插入结束标记并不在数学上阻止后续位置读取前文。

截断可能删除答案、结束标记或前置证据。对问答任务，保留完整目标和足够条件通常比机械保留前 $n$ 个词元更重要；过长样本可以分段、筛除或重新构造。训练时应报告被截断比例和空目标样本数，并禁止有效目标数为零时继续做除法。

灾难性遗忘（catastrophic forgetting）指适配新分布后已有能力显著退化。可通过较小学习率、混入通用数据、限制更新参数或增加约束缓解，但都存在任务收益与保留能力的折中。判断遗忘需要适配前后的同一组通用与领域测试，不能只凭领域验证损失下降。

\section{LoRA 的参数化、梯度与合并}

\emph{参数高效微调（parameter-efficient fine-tuning, PEFT）}只训练完整模型中的少量参数或附加模块。\emph{低秩适配（low-rank adaptation, LoRA）}冻结 $W_0\in\mathbb R^{d_o\times d_i}$，令
\begin{equation}
W=W_0+sBA,\qquad A\in\mathbb R^{r\times d_i},\quad
B\in\mathbb R^{d_o\times r},\quad s=\alpha/r.
\label{eq:lora}
\end{equation}
可训练参数量为 $r(d_i+d_o)$，低于全矩阵的条件是 $r<d_id_o/(d_i+d_o)$。低秩限制作用于更新 $\Delta W$，而不是要求原始权重 $W_0$ 低秩\cite{lora}。

设损失对有效矩阵的梯度为 $G$，则 $\nabla_B\mathcal L=sGA^\top$、$\nabla_A\mathcal L=sB^\top G$。常用初始化令 $A$ 随机、$B=0$，这样初始模型保持不变，而 $B$ 可先收到非零梯度；若二者都初始化为零，两者梯度都为零，无法开始学习。$\alpha$ 控制更新尺度，但其效果与学习率、秩和初始化相互作用。

目标模块（target modules）决定在哪些投影中加入适配器，例如查询、值、输出和 FFN；只适配部分层与全层适配表达能力不同。LoRA dropout 通常作用于适配分支输入，推理时关闭。若推理分支是固定线性更新，可合并为 $W_0+sBA$，从而消除额外分支；若基座已经量化，合并后再量化会引入新的误差，必须复测。

\begin{example}
对 $4096\times4096$ 的线性层采用 $r=16$ 的 LoRA，求可训练参数比例。
\end{example}
\begin{solve}
适配器参数为 $16(4096+4096)=131072$，完整矩阵参数为 $16777216$，比例为 $1/128\approx0.78125\%$。这只比较此层参数；训练仍需前向经过基座并保存必要激活，所以显存不会按同一比例缩小。
\end{solve}

\section{QLoRA 的存储与计算路径}

\emph{量化低秩适配（quantized low-rank adaptation, QLoRA）}以低位格式保存冻结基座，计算时按需反量化，并训练较高精度的 LoRA 参数\cite{qlora}。其输入输出仍经过完整基座，梯度可穿过线性运算流向适配器，但不更新冻结的量化基座编码。把“4位基座”理解成全部前反向均执行4位运算，是对存储与计算精度的混淆。

NF4 是根据近似正态权重分布设计的非均匀4位码本，双重量化（double quantization）进一步量化分组尺度等元数据，分页优化器（paged optimizer）利用内存管理机制缓解某些峰值。三者分别涉及权重编码、元数据压缩和内存调度，不能合并理解为一个神奇的量化公式。数据分布偏离假设、分组尺度失配或反量化开销都可能影响效果。

比较全参数微调、LoRA 与 QLoRA 时，应保持数据与有效训练预算一致，报告可训练参数、实际峰值显存、训练速度和独立任务指标。若比较的模型或数据不同，不能把全部差异归因于适配方法。量化的一般误差模型将在第~\ref{ch:quantization} 章讨论。

